\section{Introduction \& Significance}

Endometriosis is a chronic and inflammatory condition characterized by the growth of endometrial-like tissue outside the uterus\cite{Zondervan2018Endometriosis} that affects approximately 6--10\% of women of reproductive age worldwide\cite{Zondervan2018Endometriosis}. It is enigmatic in its etiopathogenesis, with no known biomarker or cure, no standards for treatment, and heterogenous symptomatic presentation\cite{Takeuchi2024Endometriosis,McKillop2019Thesis}. The diagnostic delay for endometriosis averages 7--11 years from symptom onset\cite{Agarwal2019Endometriosis,Soliman2017Endometriosis,staal2016diagnostic}. On average, patients see 4-6 clinicians before receiving a confirmation of their condition\cite{McKillop2019Thesis}.

\subsection{Communication and Management}

We argue that the personal \textit{and} clinical experience of endometriosis are shaped by a single underlying difficulty in two related ways: patients struggle to make sense of an embodied condition on their own, and struggle to convey this experience to providers in the circumscribed time permitted by a clinical encounter. This is evidenced by qualitative research that has consistently documented disconnects and misalignments in how patients and providers understand, represent, and discuss the condition\cite{Pichon2021Divided,Armour2019Endometriosis,Mardon2023EndometriosisReview,OHara2019EndometriosisReview,Denny2009Endometriosis,Young2015Endometriosis}. In light of these difficulties, we identify three allied concerns that inform the goals of this project:

\begin{enumerate}
  \item \textbf{Patients live with day-to-day uncertainty and carry complex, embodied, often inexpressed or underexpressed experience.} The day-to-day experience of endometriosis includes pain that varies in location, character, and intensity, gastrointestinal and urinary symptoms, fatigue, mood disturbance, sexual dysfunction, and functional limitation across several dimensions of health and activities of daily living\cite{McKillop2019Thesis,Urteaga2020Phenotypes}. Much of this disruptive experience is filtered through stigma, social isolation, and a history of dismissal by caregivers, which leads patients to under-report or minimize what they experience\cite{Ballard2006Endometriosis,Pichon2025Betrayal}\footnote{``\textit{It becomes its own full-time job.}''\cite{Pichon2021Divided}}.

  \item \textbf{Patients and providers operate at different temporal resolutions.} Patients live day-to-day or moment-to-moment with their condition, whereas providers think visit-to-visit. The result is that patients must compress weeks or months of embodied experience into the short and circumscribed length of of a clinical visit, whilst simultaneously recognizing and managing the emotional and cognitive labor of self-advocacy\cite{Pichon2021Divided}\footnote{``\textit{It feels like a trial, I have seven minutes as a lawyer to prove my case, to present enough evidence, to hit the particular buzzwords.}''\cite{Pichon2021Divided}.}.

  \item \textbf{Existing computational representations of the patient experience misalign with the patient's own self-assessment.} Work on Phendo, a specialized self-tracking app for endometriosis, has shown that even rigorously learned computational phenotypes match patients' own assessments of their health state only 18\% of the time, and rule-based phenotypes match only 35\% of the time\cite{Pichon2025Betrayal}. Data representations built from their patients' own self-reports fail to align with narratives of patients' embodied experiences.
\end{enumerate}

\subsection{Current Informatics Interventions}

Three classes of informatics interventions address many dimensions of this overall tension but each leaves a critical and unaddressed gap.

\begin{enumerate}
  \item \textbf{Self-tracking apps} like Phendo\cite{McKillop2019Thesis,CitizenEndoPhendo}, Bearable\cite{BearableApp}, and Visible\cite{VisibleHealth}\footnote{Phendo being an endometriosis-specific app. The others are mentioned as `general-purpose', `popular'  apps, popularity being determined by cursory examinations of both Apple App Store downloads \& searches across \texttt{MyEndoMetriosisTeam.com}} enable rich symptom data collection across many disease dimensions. They have engaged patients longitudinally\cite{Urteaga2022Engagement} supported computational phenotyping\cite{Urteaga2020Phenotypes,McKillop2019Thesis}. However, they require patients to distill the complexity of their embodied, emotional experience into structured fields (typically multi-part forms that primarily entail checkboxes, radio and other buttons, and text inputs) which yield mistranslations between the complexity and `messiness' of embodied patient experience and what an informatics intervention may utilize for downstream analysis apropos sensemaking, insight generation, or visit preparation.
  \item This under/misrepresentation suffuses \textbf{Timeline Visualizations} which do support analyst and provider decision-making processes but do not incorporate the lived experience of the patient and/or their narrative\cite{Ledesma2019HealthTimeline}.
  \item \textbf{Conversational AI for disease self-management} is emergent\cite{kocielnik2018reflection} but has not been applied to or studied in chronic or enigmatic conditions, nor has it been integrated with temporal or anatomical visualizations that support clinical communication.
\end{enumerate}

\subsection{The Gap}

We posit that \textit{there is currently no informatics intervention that maps a patient's spoken, `messy', embodied account of an enigmatic and chronic condition like endometriosis to rich and expressive self-authored, multi-temporal artifacts for reflection, sensemaking, and shared clinical decision-making.}
